\documentclass[11pt]{article}

% ---------- packages ----------
\usepackage[a4paper,margin=1in]{geometry}
\usepackage[T1]{fontenc}
\usepackage[utf8]{inputenc}
\usepackage{lmodern}
\usepackage{microtype}
\usepackage{amsmath,amssymb,amsthm,mathtools,bm}
\usepackage{booktabs}
\usepackage{longtable}
\usepackage{array}
\usepackage{graphicx}
\usepackage{subcaption}
\usepackage{multirow}
\usepackage{threeparttable}
\usepackage{adjustbox}
\usepackage{enumitem}
\usepackage{hyperref}
\usepackage{xcolor}
\usepackage{float}
\usepackage{caption}

\hypersetup{
  colorlinks=true,
  linkcolor=blue!60!black,
  citecolor=blue!60!black,
  urlcolor=blue!60!black
}

\newcommand{\R}{\mathbb{R}}
\newcommand{\E}{\mathbb{E}}
\newcommand{\Var}{\mathrm{Var}}
\newcommand{\Cov}{\mathrm{Cov}}
\newcommand{\Norm}{\mathcal{N}}

\title{Simulation Experiments for the Group-Regularized Regularized Horseshoe}
\author{}
\date{}

\begin{document}
\maketitle

\section{Simulation design}

All results reported here are based on \textbf{100 simulation replicates per
setting}. For Bayesian methods, all runs are convergence-qualified under a
common protocol using two chains, 250 warmup draws per chain, 250
post-warmup draws per chain, and automatic retry-until-convergence logic.

All experiments are based on the Gaussian linear model
\[
  y = X\beta + \varepsilon,
  \qquad
  \varepsilon \sim \Norm(0,\sigma^2 I_n).
\]
The coefficient vector \(\beta\in\R^p\) is partitioned into known groups,
\[
  \beta = (\beta_1^\top,\dots,\beta_G^\top)^\top,
  \qquad
  \beta_g \in \R^{p_g},
  \qquad
  \sum_{g=1}^G p_g = p.
\]
Within each setting, the design matrix \(X\) is sampled from a multivariate
normal distribution with block correlation structure. Each column is then
standardized. The diagonal entries of the covariance matrix are equal to one,
the within-group correlations are fixed at a setting-specific value
\(\rho_{\mathrm{within}}\), and the between-group correlations are fixed at
\(\rho_{\mathrm{between}}=0.2\). Thus, groups are genuinely correlated
internally but not independent of each other.

When a target signal-to-noise ratio is specified, \(\sigma^2\) is chosen so
that
\[
  \frac{\beta^\top \Sigma_X \beta}{\sigma^2}
\]
matches the desired population signal-to-noise ratio. When a target
coefficient of determination is specified, \(\sigma^2\) is chosen so that
\[
  \frac{\beta^\top \Sigma_X \beta}{\beta^\top \Sigma_X \beta + \sigma^2}
\]
matches the desired population \(R^2\). Train and test samples are generated
independently from the same design covariance.

Two kinds of within-group signal patterns appear repeatedly in the experiments.
In a \emph{concentrated} group, most of the group energy is carried by a small
number of coordinates. In a \emph{distributed} group, signal is spread across
many coordinates in the same group. This distinction plays the same conceptual
role that sparse-versus-dense within-group signal plays in the grouped
shrinkage literature: it determines whether a method should favor highly local
escape or broader group-level retention.

\section{Mechanism experiments}

Every mechanism setting uses five groups of size ten, so \(p=50\), together
with \(n_{\mathrm{train}}=100\) and \(n_{\mathrm{test}}=30\). The between-group
correlation is fixed at \(0.2\), while the within-group correlation is either
\(0.8\) or \(0.9\). This design keeps the problems small enough for the group
mechanism to be visually interpretable while still preserving realistic within-
group dependence.

\subsection{Visual mechanism summary}

Rather than listing all four mechanism settings separately, it is enough to keep
the two figures that show the main empirical story. The first figure shows the
group-separation pattern itself. The second figure shows that this pattern
depends on the group-specific slab component of GR-RHS rather than on a generic
increase in flexibility.

\begin{figure}[H]
\centering
\includegraphics[width=0.98\textwidth]{../outputs/history/simulation_mechanism/mechanism_main/20260426_030600_638018/figures/figure_m1_group_separation_story.png}
\caption{Group-separation mechanism under the clean ordered-signal setting.
Panel A compares posterior mean \(\kappa_g\) values for null and signal groups
across 100 replicates. Panel B orders the five groups by true group strength
\((0,0,1.5,4,10)\) and shows the corresponding mean posterior retention. Null
groups remain near zero, while weak, medium, and strong signal groups receive
progressively larger retention.}
\label{fig:m1-group-separation}
\end{figure}

Figure~\ref{fig:m1-group-separation} is the cleanest visualization of the
retention-ceiling interpretation of \(\kappa_g\). Across the 100 replicates,
GR-RHS achieves perfect group AUROC and a mean \texttt{kappa\_gap} of about
\(0.29\). More importantly, the figure shows the mechanism in directly
interpretable form: null groups are held close to zero, and the posterior
retention rises monotonically with true group strength.

\begin{figure}[H]
\centering
\includegraphics[width=0.98\textwidth]{../outputs/history/simulation_mechanism/mechanism_main/20260426_030600_638018/figures/figure_m4_mechanism_attribution.png}
\caption{Mechanism attribution by ablation. Panel A reports the change in
\texttt{kappa\_gap} relative to the full GR-RHS model. Panel B reports the
effective \(\tau_0\) calibration relative to the oracle baseline. Oracle
\(\tau_0\) and fixed \(10\times \tau_0\) preserve group separation, removing
local scales weakens it, and forcing shared or absent \(\kappa\) eliminates it
almost completely.}
\label{fig:m4-ablation}
\end{figure}

Figure~\ref{fig:m4-ablation} explains where that pattern comes from. Changing
the global calibration alone does not materially remove group separation: the
oracle-\(\tau_0\) and fixed-\(10\times\tau_0\) variants stay close to the full
model in Panel A. Removing local scales weakens the effect, but the decisive
result is that shared \(\kappa\) and no \(\kappa\) both collapse the
\texttt{kappa\_gap} by about \(0.41\), essentially wiping out the mechanism.
This is the core empirical reason to interpret GR-RHS as a grouped slab model:
the visible separation is created mainly by the group-specific slab ceiling,
not by generic overparameterization.

\section{Main benchmark experiments}

\subsection{Generative model}

The data-generating mechanism for the benchmark experiments is the Gaussian
linear regression model
\[
  y = X\beta + \varepsilon,
  \qquad
  \varepsilon \sim \Norm(0,\sigma^2 I_n),
\]
where \(X\) is sampled from a multivariate normal distribution with mean zero
and grouped block correlation matrix \(\Sigma_X\). The covariance matrix is
constructed so that each regressor has unit variance, the within-group
correlations are exchangeable, and the pairwise correlations across groups are
equal to \(0.2\). For the medium-correlation benchmark settings,
\(\rho_{\mathrm{within}}=0.6\), while for the high-correlation benchmark
settings, \(\rho_{\mathrm{within}}=0.8\). In all settings, the residual error
variance is chosen so that
\[
  \frac{\beta^\top \Sigma_X \beta}{\beta^\top \Sigma_X \beta + \sigma^2}=0.7.
\]

The benchmark settings in this paper are random-coefficient grouped regression
settings rather than fixed-coefficient settings. In each replicate, the true
coefficient vector is generated from a family-specific grouped signal blueprint.
The benchmark contains a low-dimensional regime only, with
\[
  n_{\mathrm{train}}=500,\qquad n_{\mathrm{test}}=100,\qquad p=50,
\]
except that the grouping pattern varies across settings. The methods compared
are GR-RHS, RHS, Grouped Horseshoe+, GIGG-MMLE, Lasso with cross-validation,
and OLS.

Three random signal families are used. In the \emph{classical reference}
family, the active groups share a common mild within-group blueprint. The
active support fraction is drawn uniformly from \((0.6,0.8)\), and the
concentration parameter controlling the internal allocation of group energy is
drawn uniformly from \((3,6)\). In the \emph{single-mode heterogeneous}
family, the active groups still share a common blueprint, but the support
fraction is drawn from \((0.4,0.8)\) and the concentration parameter is drawn
from \((0.8,2)\), producing a more heterogeneous within-group signal profile.
In the \emph{multimode heterogeneous} family, the active groups no longer share
the same blueprint. Instead, each active group draws its own support fraction
from \((0.2,0.8)\) and its own concentration parameter from a log-uniform
distribution on \((0.2,6)\). Additional acceptance criteria are imposed so
that the active groups do not collapse to the same signal mode; in particular,
the generated active groups must display sufficient separation either in
support fraction or in concentration.

Conditional on the family draw, signal energy is allocated across the active
groups by a Dirichlet distribution. Within each active group, the support set is
sampled uniformly over coordinates and the active coefficients are assigned a
common sign within group. The resulting coefficient vector therefore varies from
replicate to replicate, but always preserves the grouped signal architecture of
the benchmark family. The purpose of this random-coefficient benchmark is to
measure average performance across many grouped sparse configurations rather
than only on a small number of hand-tuned coefficient vectors.

\subsection{Benchmark settings}

The six benchmark settings are summarized in
Table~\ref{tab:benchmark-settings}. All settings use
\(n_{\mathrm{train}}=500\), \(n_{\mathrm{test}}=100\), target \(R^2=0.7\), and
\(\rho_{\mathrm{between}}=0.2\).

\begin{table}[H]
\centering
\caption{Benchmark setting definitions.}
\label{tab:benchmark-settings}
\begin{tabular}{cllccc}
\toprule
ID & Family & Group sizes & Active groups & \(\rho_{\mathrm{within}}\) & Main role \\
\midrule
1 & classical reference & \((10,10,10,10,10)\) & \(3/5\) & 0.6 & anchor \\
2 & classical reference & \((10,10,10,10,10)\) & \(3/5\) & 0.8 & higher correlation \\
3 & single-mode heterogeneous & \((10,10,10,10,10)\) & \(3/5\) & 0.6 & equal-size heterogeneous \\
4 & single-mode heterogeneous & \((30,10,5,3,2)\) & \(3/5\) & 0.6 & unequal groups \\
5 & multimode heterogeneous & \((10,10,10,10,10)\) & \(3/5\) & 0.8 & equal-size multimode \\
6 & multimode heterogeneous & \((25,25)\) & \(2/2\) & 0.8 & large-group multimode \\
\bottomrule
\end{tabular}
\end{table}

Settings 1 and 2 are classical equal-group benchmarks and differ only by the
strength of within-group correlation. Setting 3 keeps the equal-group geometry
but replaces the classical signal family with a single-mode heterogeneous
family. Setting 4 uses the same signal family as Setting 3 but changes the
group sizes to \((30,10,5,3,2)\), creating a benchmark with substantial size
imbalance across groups. Setting 5 is the equal-size multimode benchmark, in
which the active groups are allowed to differ substantially in their internal
signal structure. Setting 6 is the large-group multimode benchmark with two
groups of size 25. These six settings move from relatively standard grouped
regression designs toward settings in which a single global slab width is least
well aligned with the true grouped signal geometry.

\section{High-dimensional sampler design and literature sources}

This section documents the sources behind the high-dimensional GR-RHS sampling
design and clarifies which components are inherited from existing Bayesian
shrinkage and MCMC ideas and which components are specific to our model. The
high-dimensional sampler used for the \(p>n\) extension is called
\textbf{GR-RHS-CaP}, short for \emph{collapsed-and-profile GR-RHS}. Its purpose
is not to replace the staged Gibbs sampler in the low-dimensional Gaussian
experiments. Rather, it provides a more scalable inference path for
high-dimensional grouped regression, where coefficient-wise local-scale updates
become the dominant computational cost.

\subsection{Modeling ingredients inherited from the shrinkage literature}

The global-local shrinkage principle used by GR-RHS follows the horseshoe
line of work. Carvalho, Polson, and Scott (2010), ``The horseshoe estimator for
sparse signals,'' introduced the horseshoe estimator as a prior that strongly
shrinks small coefficients while allowing large coefficients to escape through
local scales. Polson and Scott (2010), ``Shrink globally, act locally: sparse
Bayesian regularization and prediction,'' gives the same idea in the language
of a global scale controlling overall sparsity and local scales protecting
large signals. In GR-RHS, the global scale \(\tau\) and coefficient-level local
scales \(\lambda_j\) inherit this global-local logic.

The finite-slab part of the model is motivated by the regularized horseshoe
literature. Piironen and Vehtari (2017), especially their work on sparsity
information and regularization in the horseshoe prior, introduced the practical
idea that the horseshoe should be regularized by a slab-like upper scale so
that very large coefficients do not become effectively unregularized. The
\texttt{rstanarm} implementation of the regularized horseshoe follows this
same finite-slab design. GR-RHS keeps this regularized-horseshoe motivation but
changes the unit at which the slab ceiling is learned: instead of using only a
single coefficient-blind slab scale, GR-RHS introduces group-specific retention
parameters \(\kappa_g\).

Grouped sparsity as a statistical target is connected to the group-lasso and
sparse-group-lasso literature. Yuan and Lin (2006), ``Model selection and
estimation in regression with grouped variables,'' established the importance
of treating predefined groups as units of regularization. Meier, van de Geer,
and Buhlmann (2008) extended grouped regularization ideas to logistic
regression, and Simon, Friedman, Hastie, and Tibshirani (2013) developed the
sparse-group lasso, combining group-level and within-group sparsity. GR-RHS is
Bayesian rather than penalized-likelihood based, but the simulation design uses
the same scientific distinction: a method may need to retain or suppress whole
groups while still allowing sparse or heterogeneous behavior within active
groups.

The grouped horseshoe baseline is related to the grouped global-local
shrinkage literature, while the Horseshoe+ baseline follows Bhadra, Datta,
Polson, and Willard (2017), ``The Horseshoe+ estimator of ultra-sparse
signals.'' In our experiments, these baselines help separate the benefit of
generic group-aware shrinkage from the specific GR-RHS mechanism, namely a
group-specific slab-retention ceiling.

\subsection{Sampling ingredients inherited from the MCMC literature}

The key high-dimensional Gaussian sampling identity comes from Bhattacharya,
Chakraborty, and Mallick (2016), ``Fast sampling with Gaussian scale mixture
priors in high-dimensional regression.'' They showed how to draw regression
coefficients from a Gaussian conditional posterior under scale-mixture priors
using a Woodbury/Bhattacharya sampler with complexity favorable when
\(n \ll p\). GR-RHS-CaP uses this idea for the conditional draw
\[
  \beta \mid \theta, y,
  \qquad
  \theta = (\sigma,\tau,\lambda,\kappa),
\]
after the hyperparameters have been sampled. In the implementation, this is
the same computational principle used by the repository's
\texttt{beta\_sample\_woodbury} routine.

The dynamic Hamiltonian Monte Carlo component follows Hoffman and Gelman
(2014), ``The No-U-Turn Sampler: adaptively setting path lengths in
Hamiltonian Monte Carlo.'' GR-RHS-CaP does not introduce a new HMC algorithm.
It uses NUTS as implemented in NumPyro, but changes the posterior geometry that
NUTS sees by marginalizing \(\beta\) out of the main sampling state.

The collapsed-sampling idea is a standard Rao-Blackwellized Bayesian
computation strategy: analytically integrate a conditionally Gaussian block
when possible, sample the lower-dimensional marginal posterior, and then
reconstruct the integrated block from its exact conditional posterior. In
GR-RHS-CaP, the marginal likelihood is
\[
  y \mid \theta \sim
  \Norm\!\left(0,\, X V(\theta) X^\top + \sigma^2 I_n\right),
\]
where \(V(\theta)\) is the diagonal prior covariance implied by the
regularized horseshoe hierarchy. This is the part that makes the high-
dimensional sampler substantially different from a joint NUTS sampler over
\(\beta\), \(\lambda\), \(\tau\), and \(\kappa\).

\subsection{What is specific to GR-RHS-CaP}

The first model-specific contribution is the \emph{group-specific slab
retention} parameter \(\kappa_g\). In ordinary RHS, the regularized slab is
coefficient-blind. In GR-RHS, each coefficient belongs to a group \(g(j)\), and
the regularized prior variance uses the group-level retention parameter
\(\kappa_{g(j)}\). This makes \(\kappa_g\) interpretable as a group-level
ceiling on how much signal can escape global shrinkage. The mechanism
experiments in this document are designed around exactly this interpretation:
they evaluate whether signal groups receive larger posterior retention than
null groups and whether this separation disappears when the group gate is
removed.

The second model-specific contribution is the \emph{profile mode} used for
high-dimensional experiments. When \(\texttt{use\_local\_scale=false}\), the
coefficient-level local scales are profiled out by setting \(\lambda_j=1\). The
main NUTS state then contains essentially \((\sigma,\tau,\kappa_1,\dots,
\kappa_G)\), so the dimension is approximately \(G+2\) rather than \(p+G+2\).
This is the reason the high-dimensional configuration uses
\[
  \texttt{sampler\_backend = collapsed\_profile},
  \qquad
  \texttt{use\_local\_scale = false}.
\]
The design is intentionally conservative: in very high dimensions, the first
goal is to learn stable group-level retention before reintroducing thousands of
coefficient-specific local scales.

The third model-specific contribution is the two-resolution interpretation of
GR-RHS inference. The low-dimensional and moderate-dimensional Gaussian
experiments use the staged Gibbs sampler because full coefficient-level
posterior sampling is still feasible there. The high-dimensional experiments
use GR-RHS-CaP because the coefficient block can be reconstructed after
sampling the lower-dimensional grouped hyperparameter posterior. Thus the
repository now uses two inference regimes for the same statistical prior:
staged Gibbs for detailed low-dimensional posterior exploration and
collapsed-profile NUTS for scalable high-dimensional grouped screening.

\subsection{What is original in the experimental design}

The random-coefficient benchmark design is our own validation layer around the
GR-RHS mechanism. Rather than testing only fixed coefficient vectors, each
replicate draws active groups, within-group support fractions, energy shares,
and concentration parameters. This creates benchmark families in which the
same method must work across many grouped sparse configurations. The
single-mode and multimode families are especially important because they test
whether a group-aware slab is useful when active groups have heterogeneous
internal signal allocation.

The high-dimensional extension is also our own adaptation to the GR-RHS
architecture. The borrowed computational ideas are clear: global-local
shrinkage from the horseshoe literature, finite-slab regularization from the
regularized horseshoe, fast Gaussian conditional sampling from Bhattacharya,
Chakraborty, and Mallick, and NUTS from Hoffman and Gelman. What is specific
to this project is the way those components are assembled around
\(\kappa_g\): the finite slab is made group-specific, the high-dimensional
sampler collapses \(\beta\), the profile mode reduces the NUTS state to
group-level retention parameters, and the experiments explicitly test whether
that group-level retention is learned in the regimes where grouped structure is
informative.

\subsection{Reference checklist for the sampler and design}

For clarity, the main external sources used in the design are:
\begin{itemize}[leftmargin=*]
  \item Carvalho, Polson, and Scott (2010), ``The horseshoe estimator for
  sparse signals'': global-local shrinkage and local escape.
  \item Polson and Scott (2010), ``Shrink globally, act locally: sparse
  Bayesian regularization and prediction'': global shrinkage with local
  adaptivity as the organizing principle.
  \item Piironen and Vehtari (2017), ``Sparsity information and regularization
  in the horseshoe and other shrinkage priors'': finite-slab regularization and
  practical sparsity calibration for the regularized horseshoe.
  \item Yuan and Lin (2006), ``Model selection and estimation in regression
  with grouped variables'': group-level regularization as a statistical target.
  \item Meier, van de Geer, and Buhlmann (2008), ``The group lasso for logistic
  regression'': grouped regularization beyond Gaussian least squares.
  \item Simon, Friedman, Hastie, and Tibshirani (2013), ``A sparse-group
  lasso'': simultaneous group-level and within-group sparsity.
  \item Bhadra, Datta, Polson, and Willard (2017), ``The Horseshoe+ estimator
  of ultra-sparse signals'': the ultra-sparse Horseshoe+ baseline family.
  \item Bhattacharya, Chakraborty, and Mallick (2016), ``Fast sampling with
  Gaussian scale mixture priors in high-dimensional regression'': Woodbury
  conditional Gaussian sampling for \(n \ll p\).
  \item Hoffman and Gelman (2014), ``The No-U-Turn Sampler: adaptively setting
  path lengths in Hamiltonian Monte Carlo'': the NUTS engine used for the
  collapsed hyperparameter posterior.
\end{itemize}

\section{Evaluation criteria}

The main evaluation metrics are test-set overall mean-squared error, signal
mean-squared error, null mean-squared error, test log predictive density, and
95\% posterior interval coverage for Bayesian methods. The overall MSE measures
the average squared error across all coefficients, while signal MSE and null MSE
separate the active and inactive coordinates. This distinction is useful because
a method can attain a low overall error either by estimating active
coefficients well or by aggressively shrinking null coefficients, and these are
not the same behavior.

For the mechanism experiments, additional group-level summaries are recorded.
The most important are the group AUROC and the quantity called
\texttt{kappa\_gap}, which measures the average difference between signal-group
and null-group posterior retention. In the context of GR-RHS, these summaries
are natural because the theory interprets \(\kappa_g\) as a group-level
retention ceiling.

\section{Benchmark results}

The benchmark results are reported in two complementary forms. First,
Table~\ref{tab:benchmark-full} gives the full six-setting by six-method summary
for the main error criteria. For each method and setting, the table reports the
replicate mean and replicate variance of overall MSE, signal MSE, and null MSE.
It also reports the replicate mean and replicate variance of test LPD. For
Bayesian methods, the table additionally reports empirical mean coverage and
mean interval width for nominal 95\% intervals. Second,
Figures~\ref{fig:bench-s1}--\ref{fig:bench-s6} provide setting-specific plots
for the Bayesian methods, combining MSE and coverage in a single display.

\begin{table}[H]
\centering
\caption{Method numbering used in the benchmark result displays.}
\label{tab:method-index}
\begin{tabular}{cl}
\toprule
Code & Method \\
\midrule
M1 & GR-RHS \\
M2 & RHS \\
M3 & Grouped Horseshoe+ \\
M4 & GIGG-MMLE \\
M5 & Lasso (CV) \\
M6 & OLS \\
\bottomrule
\end{tabular}
\end{table}

\begin{table}[p]
\centering
\caption{Full benchmark summary for the six settings and six methods. Each MSE
entry is reported in the form mean \(\pm\) variance across 100 replicates.
LPD is also reported as mean \(\pm\) variance. Coverage and interval width are
reported only for Bayesian methods.}
\label{tab:benchmark-full}
\begin{adjustbox}{max width=\textwidth}
\begin{tabular}{llccccccc}
\toprule
Setting & Method & Overall MSE & Signal MSE & Null MSE & LPD & Interval Width & Coverage \\
\midrule
Setting 1 & GR-RHS & 0.004691 $\pm$ 4.67e-06 & \textbf{0.008181 $\pm$ 1.59e-05} & 0.002139 $\pm$ 2.23e-06 & \textbf{-1.716177 $\pm$ 2.57e-02} & 0.2866 & \textbf{0.9732} \\
Setting 1 & RHS & 0.005491 $\pm$ 5.74e-06 & 0.010106 $\pm$ 2.15e-05 & 0.002126 $\pm$ 2.37e-06 & -1.722641 $\pm$ 2.58e-02 & 0.2903 & 0.9670 \\
Setting 1 & Grouped Horseshoe+ & 0.004895 $\pm$ 4.79e-06 & 0.009159 $\pm$ 1.73e-05 & 0.001776 $\pm$ 2.13e-06 & -1.717428 $\pm$ 2.55e-02 & 0.2808 & 0.9700 \\
Setting 1 & GIGG-MMLE & 0.004920 $\pm$ 5.74e-06 & 0.009959 $\pm$ 1.99e-05 & \textbf{0.001210 $\pm$ 2.02e-06} & -1.718012 $\pm$ 2.67e-02 & 0.0583 & 0.2792 \\
Setting 1 & Lasso (CV) & \textbf{0.004513 $\pm$ 4.54e-06} & 0.008552 $\pm$ 1.50e-05 & 0.001553 $\pm$ 2.33e-06 & -1.721452 $\pm$ 2.68e-02 & -- & -- \\
Setting 1 & OLS & 0.009210 $\pm$ 1.76e-05 & 0.009257 $\pm$ 2.25e-05 & 0.009152 $\pm$ 2.03e-05 & -1.748034 $\pm$ 2.84e-02 & -- & -- \\
\midrule
Setting 2 & GR-RHS & 0.009878 $\pm$ 1.45e-05 & \textbf{0.016870 $\pm$ 4.37e-05} & 0.004792 $\pm$ 1.41e-05 & \textbf{-1.867236 $\pm$ 2.22e-02} & 0.4391 & \textbf{0.9772} \\
Setting 2 & RHS & 0.011565 $\pm$ 1.80e-05 & 0.021111 $\pm$ 5.81e-05 & 0.004646 $\pm$ 1.61e-05 & -1.872078 $\pm$ 2.35e-02 & 0.4200 & 0.9672 \\
Setting 2 & Grouped Horseshoe+ & 0.010440 $\pm$ 1.57e-05 & 0.018682 $\pm$ 4.74e-05 & 0.004446 $\pm$ 1.49e-05 & -1.868290 $\pm$ 2.23e-02 & 0.4357 & 0.9758 \\
Setting 2 & GIGG-MMLE & 0.011585 $\pm$ 2.84e-05 & 0.022153 $\pm$ 5.63e-05 & 0.003831 $\pm$ 3.77e-05 & -1.871365 $\pm$ 2.46e-02 & 0.0842 & 0.2830 \\
Setting 2 & Lasso (CV) & \textbf{0.009850 $\pm$ 1.37e-05} & 0.018215 $\pm$ 3.82e-05 & \textbf{0.003736 $\pm$ 1.24e-05} & -1.871942 $\pm$ 2.46e-02 & -- & -- \\
Setting 2 & OLS & 0.025247 $\pm$ 9.13e-05 & 0.025268 $\pm$ 1.23e-04 & 0.025236 $\pm$ 1.28e-04 & -1.913207 $\pm$ 2.48e-02 & -- & -- \\
\midrule
Setting 3 & GR-RHS & 0.003697 $\pm$ 3.15e-06 & \textbf{0.006781 $\pm$ 1.02e-05} & 0.001813 $\pm$ 1.35e-06 & -1.644194 $\pm$ 3.85e-02 & 0.2618 & \textbf{0.9792} \\
Setting 3 & RHS & 0.004046 $\pm$ 4.19e-06 & 0.007962 $\pm$ 1.38e-05 & 0.001664 $\pm$ 1.38e-06 & -1.647789 $\pm$ 3.92e-02 & 0.2588 & 0.9768 \\
Setting 3 & Grouped Horseshoe+ & 0.003701 $\pm$ 3.37e-06 & 0.007241 $\pm$ 1.15e-05 & 0.001516 $\pm$ 1.12e-06 & -1.644024 $\pm$ 3.87e-02 & 0.2545 & 0.9780 \\
Setting 3 & GIGG-MMLE & 0.003573 $\pm$ 5.04e-06 & 0.007679 $\pm$ 1.41e-05 & \textbf{0.001019 $\pm$ 2.18e-06} & \textbf{-1.642351 $\pm$ 4.02e-02} & 0.0544 & 0.3004 \\
Setting 3 & Lasso (CV) & \textbf{0.003406 $\pm$ 3.26e-06} & 0.006838 $\pm$ 9.99e-06 & 0.001302 $\pm$ 1.66e-06 & -1.646136 $\pm$ 4.01e-02 & -- & -- \\
Setting 3 & OLS & 0.007727 $\pm$ 1.16e-05 & 0.007239 $\pm$ 1.46e-05 & 0.007954 $\pm$ 1.44e-05 & -1.684711 $\pm$ 4.07e-02 & -- & -- \\
\midrule
Setting 4 & GR-RHS & \textbf{0.005155 $\pm$ 4.04e-06} & \textbf{0.006648 $\pm$ 6.80e-06} & 0.003386 $\pm$ 4.01e-06 & \textbf{-1.765877 $\pm$ 3.16e-02} & 0.3113 & \textbf{0.9754} \\
Setting 4 & RHS & 0.005583 $\pm$ 5.26e-06 & 0.007860 $\pm$ 9.48e-06 & 0.002785 $\pm$ 3.79e-06 & -1.768928 $\pm$ 3.19e-02 & 0.2995 & 0.9664 \\
Setting 4 & Grouped Horseshoe+ & 0.005492 $\pm$ 5.06e-06 & 0.007575 $\pm$ 9.00e-06 & 0.002924 $\pm$ 3.95e-06 & -1.766971 $\pm$ 3.17e-02 & 0.3078 & 0.9706 \\
Setting 4 & GIGG-MMLE & 0.006229 $\pm$ 1.14e-05 & 0.008585 $\pm$ 1.43e-05 & 0.003308 $\pm$ 1.10e-05 & -1.773102 $\pm$ 3.41e-02 & 0.0798 & 0.3530 \\
Setting 4 & Lasso (CV) & 0.005321 $\pm$ 5.09e-06 & 0.007494 $\pm$ 8.27e-06 & \textbf{0.002647 $\pm$ 4.60e-06} & -1.773274 $\pm$ 3.32e-02 & -- & -- \\
Setting 4 & OLS & 0.010139 $\pm$ 1.92e-05 & 0.010209 $\pm$ 2.14e-05 & 0.010112 $\pm$ 2.81e-05 & -1.801798 $\pm$ 3.21e-02 & -- & -- \\
\midrule
Setting 5 & GR-RHS & 0.005347 $\pm$ 7.71e-06 & 0.012536 $\pm$ 4.35e-05 & 0.002451 $\pm$ 2.51e-06 & -1.587631 $\pm$ 3.58e-02 & 0.3234 & 0.9772 \\
Setting 5 & RHS & 0.005695 $\pm$ 9.77e-06 & 0.014901 $\pm$ 5.99e-05 & 0.001971 $\pm$ 2.54e-06 & -1.590471 $\pm$ 3.72e-02 & 0.2946 & 0.9738 \\
Setting 5 & Grouped Horseshoe+ & 0.005405 $\pm$ 8.00e-06 & 0.013070 $\pm$ 4.62e-05 & 0.002289 $\pm$ 2.36e-06 & -1.587934 $\pm$ 3.57e-02 & 0.3198 & \textbf{0.9774} \\
Setting 5 & GIGG-MMLE & 0.004939 $\pm$ 8.90e-06 & 0.013500 $\pm$ 5.21e-05 & \textbf{0.001430 $\pm$ 2.24e-06} & \textbf{-1.585691 $\pm$ 3.70e-02} & 0.0663 & 0.3590 \\
Setting 5 & Lasso (CV) & \textbf{0.004801 $\pm$ 7.11e-06} & \textbf{0.012219 $\pm$ 3.83e-05} & 0.001792 $\pm$ 2.41e-06 & -1.588458 $\pm$ 3.75e-02 & -- & -- \\
Setting 5 & OLS & 0.013939 $\pm$ 3.98e-05 & 0.014418 $\pm$ 5.50e-05 & 0.013748 $\pm$ 4.68e-05 & -1.630197 $\pm$ 4.01e-02 & -- & -- \\
\midrule
Setting 6 & GR-RHS & \textbf{0.011911 $\pm$ 1.30e-05} & \textbf{0.013789 $\pm$ 2.18e-05} & 0.010161 $\pm$ 1.63e-05 & -2.025176 $\pm$ 4.05e-02 & 0.5511 & \textbf{0.9842} \\
Setting 6 & RHS & 0.013383 $\pm$ 3.20e-05 & 0.020042 $\pm$ 6.80e-05 & \textbf{0.006450 $\pm$ 1.68e-05} & \textbf{-2.021275 $\pm$ 4.40e-02} & 0.4901 & 0.9734 \\
Setting 6 & Grouped Horseshoe+ & 0.012939 $\pm$ 2.85e-05 & 0.018699 $\pm$ 5.97e-05 & 0.006921 $\pm$ 1.57e-05 & -2.024206 $\pm$ 4.21e-02 & 0.5375 & 0.9820 \\
Setting 6 & GIGG-MMLE & 0.020460 $\pm$ 1.30e-04 & 0.026135 $\pm$ 1.60e-04 & 0.014624 $\pm$ 1.47e-04 & -2.031358 $\pm$ 4.75e-02 & 0.1210 & 0.3266 \\
Setting 6 & Lasso (CV) & 0.013569 $\pm$ 3.29e-05 & 0.019816 $\pm$ 5.65e-05 & 0.006949 $\pm$ 2.42e-05 & -2.022468 $\pm$ 4.51e-02 & -- & -- \\
Setting 6 & OLS & 0.038278 $\pm$ 2.62e-04 & 0.038405 $\pm$ 3.78e-04 & 0.037717 $\pm$ 3.08e-04 & -2.063961 $\pm$ 4.72e-02 & -- & -- \\
\bottomrule
\end{tabular}
\end{adjustbox}
\end{table}

Table~\ref{tab:benchmark-full} should be read as a joint comparison of point
estimation quality, predictive performance, and uncertainty calibration. The
overall MSE column measures aggregate coefficient recovery across all
coordinates, whereas the signal and null MSE columns separate performance on
truly active and truly inactive coefficients. This separation is important in
grouped sparse regression because a low overall MSE can arise either from good
recovery of the active coefficients or from aggressive shrinkage of the null
coefficients, and these represent substantively different behaviors. The LPD
column adds a predictive criterion based on the test distribution rather than
only on coefficient error. Finally, the interval-width and coverage columns are
meant to be interpreted together: narrower intervals are useful only when they
retain coverage close to the nominal 95\% level.

Several broad patterns are immediately visible. First, GR-RHS is the strongest
method in signal recovery. It attains the best signal MSE in five of the six
benchmark settings and is especially strong in the settings with more
informative grouped structure. This is the most direct evidence that the method
is helping for the right reason. Its main gain is not merely stronger shrinkage
of null coefficients; rather, it more consistently preserves and estimates the
nonzero part of the signal. The fact that this advantage appears across both
equal-group and heterogeneous-group benchmarks suggests that the grouped slab
layer is responding to signal geometry rather than to a single idiosyncratic
design.

Second, the overall-MSE ranking is more scope-dependent. Lasso is the best
overall-MSE method in Settings 1, 2, 3, and 5, while GR-RHS is best in
Settings 4 and 6. This pattern is substantively meaningful. In the more
classical or only mildly heterogeneous settings, GR-RHS behaves like a strong
Bayesian competitor but does not force an artificial advantage over well-tuned
classical regularization. Once the grouped structure becomes more informative,
however, the group-specific slab begins to translate into an overall-error
advantage. Setting 4, with substantial group-size imbalance, and Setting 6,
with large multimode groups, are the clearest examples of the regime for which
GR-RHS was designed.

Third, the null-MSE column shows that some methods can look attractive by
shrinking inactive coordinates more aggressively. In particular, GIGG-MMLE and
Lasso often have very favorable null-MSE values. This is useful information,
but it should not be overinterpreted in isolation. Null MSE rewards strong
suppression of noise coordinates, whereas signal MSE reveals whether the same
method is preserving true effects. The benchmark therefore argues against any
one-dimensional reading of sparsity performance. A method can be excellent at
driving null coefficients toward zero and still be inferior at recovering the
active part of the signal.

The LPD column provides a predictive perspective that is related to, but not
identical with, coefficient error. GR-RHS has the best LPD in Settings 1, 2,
and 4, while GIGG-MMLE is best in Settings 3 and 5 and RHS is best in Setting
6. This shows that predictive density does not mechanically track MSE. In some
settings, a method can obtain competitive predictive scores even when its
coefficient recovery is less favorable. For this reason, the benchmark is more
informative when LPD is read together with the uncertainty summaries. A method
that gains predictive sharpness at the expense of calibration should be treated
differently from one that improves predictive fit while maintaining reliable
uncertainty quantification.

This is exactly where the interval-width and coverage columns become important.
GR-RHS, RHS, and Grouped Horseshoe+ all maintain empirical coverage close to
nominal across the six settings, with values roughly between 0.97 and 0.98.
Their average interval widths are also of similar magnitude, indicating a
comparable sharpness-calibration tradeoff. By contrast, GIGG-MMLE is much
narrower in every setting, but its coverage collapses to roughly 0.28--0.36.
That pattern should not be interpreted as superior efficiency. It is instead a
classic sign of undercalibrated uncertainty: the intervals are narrow because
they are overconfident, not because they are better matched to the data.

Taken together, the table supports a scope-conditional interpretation of the
benchmark. GR-RHS is not a universally dominant procedure in every low-
dimensional grouped regression problem, and the table does not suggest that it
should be. Rather, it is a method that becomes increasingly advantageous when
the grouped structure is genuinely informative for shrinkage. In those regimes,
it improves signal recovery, can improve overall error, often improves
predictive fit, and does so without sacrificing calibration. That combination
of strong signal MSE and reliable coverage is the central empirical argument in
favor of the GR-RHS prior.

The overall-MSE winners are:
\[
  \text{S1: M5},\quad
  \text{S2: M5},\quad
  \text{S3: M5},\quad
  \text{S4: M1},\quad
  \text{S5: M5},\quad
  \text{S6: M1}.
\]
Thus GR-RHS is the overall winner in \(2/6\) settings, namely S4 and S6. The
signal-MSE winners are:
\[
  \text{S1: M1},\quad
  \text{S2: M1},\quad
  \text{S3: M1},\quad
  \text{S4: M1},\quad
  \text{S5: M5},\quad
  \text{S6: M1}.
\]
Hence GR-RHS is the signal-MSE winner in \(5/6\) settings.

For coverage, the Bayesian methods satisfy the ordering
\[
  \text{GR-RHS} \approx \text{RHS} \approx \text{GHS+} \gg \text{GIGG-MMLE}
\]
in every setting. Numerically, GR-RHS coverage ranges from \(0.9732\) to
\(0.9842\), RHS from \(0.9664\) to \(0.9768\), Grouped Horseshoe+ from
\(0.9700\) to \(0.9820\), and GIGG-MMLE from \(0.2792\) to \(0.3590\).

For interval width, the same table shows the usual calibration-sharpness tradeoff.
GR-RHS, RHS, and Grouped Horseshoe+ have broadly similar average widths, while
GIGG-MMLE is much narrower in every setting. The narrowness of GIGG-MMLE,
however, comes together with severe undercoverage rather than better-calibrated
uncertainty.

The highest replicate-level MSE variances occur in Setting 6. For example,
GIGG-MMLE reaches \(\Var(\mathrm{MSE}_{\mathrm{overall}})=1.30\times 10^{-4}\)
and OLS reaches \(2.62\times 10^{-4}\), both substantially larger than in the
other settings.

\begin{figure}[p]
\centering
\includegraphics[width=0.95\textwidth]{generated_figures/setting_1_classical_equal_medium_bayes_mse_coverage.png}
\caption{Bayesian-method comparison for Setting 1. The left panel reports the
mean overall, signal, and null MSE, with one-standard-deviation error bars
across 100 replicates. The right panel reports empirical 95\% coverage for the
Bayesian methods.}
\label{fig:bench-s1}
\end{figure}

\begin{figure}[p]
\centering
\includegraphics[width=0.95\textwidth]{generated_figures/setting_2_classical_equal_high_bayes_mse_coverage.png}
\caption{Bayesian-method comparison for Setting 2. The design is the same as
Setting 1 except that within-group correlation increases from \(0.6\) to
\(0.8\).}
\label{fig:bench-s2}
\end{figure}

\begin{figure}[p]
\centering
\includegraphics[width=0.95\textwidth]{generated_figures/setting_3_single_mode_equal_bayes_mse_coverage.png}
\caption{Bayesian-method comparison for Setting 3, the equal-size single-mode
heterogeneous benchmark.}
\label{fig:bench-s3}
\end{figure}

\begin{figure}[p]
\centering
\includegraphics[width=0.95\textwidth]{generated_figures/setting_4_single_mode_unequal_bayes_mse_coverage.png}
\caption{Bayesian-method comparison for Setting 4, the unequal-group
single-mode heterogeneous benchmark. This is one of the first settings in which
the grouped layer becomes clearly useful in overall MSE.}
\label{fig:bench-s4}
\end{figure}

\begin{figure}[p]
\centering
\includegraphics[width=0.95\textwidth]{generated_figures/setting_5_multimode_equal_bayes_mse_coverage.png}
\caption{Bayesian-method comparison for Setting 5, the equal-size multimode
heterogeneous showcase.}
\label{fig:bench-s5}
\end{figure}

\begin{figure}[p]
\centering
\includegraphics[width=0.95\textwidth]{generated_figures/setting_6_multimode_large_groups_bayes_mse_coverage.png}
\caption{Bayesian-method comparison for Setting 6, the large-group multimode
benchmark. This is the most favorable benchmark regime for a grouped slab
model, and it is also the regime with the clearest overall win for GR-RHS.}
\label{fig:bench-s6}
\end{figure}

Figures~\ref{fig:bench-s1}--\ref{fig:bench-s6} present the same information
from the viewpoint of the Bayesian competitors only. This is useful because the
coverage comparison is most meaningful within the subset of methods that
actually produce posterior intervals. Across all six settings, GR-RHS maintains
coverage close to nominal. RHS and Grouped Horseshoe+ do the same, although
their error levels are not always as favorable in the grouped settings. By
contrast, GIGG-MMLE often attains competitive point error but exhibits very low
coverage, with empirical values around \(0.28\) to \(0.36\), far below the
nominal 0.95 line shown in the right panels.

\section{Real data experiment}

\subsection{Dataset source and response}

To complement the synthetic benchmarks, we also evaluate the methods on a real
DNA methylation age dataset derived from GEO accession \texttt{GSE40279}. This
study contains human whole-blood methylation measurements profiled on the
Illumina HumanMethylation450 array, together with chronological age for each
sample. We use the processed beta-value matrix released with the series,
together with the GEO sample key, the series-level metadata, and the Illumina
450k platform annotation. The response variable is chronological age in years,
parsed from the GEO sample metadata. After aligning the sample order between
the age annotations and the beta-value matrix, the final dataset contains
\(n=656\) subjects, with age ranging from \(19\) to \(101\) years and mean age
approximately \(64.0\) years.

\subsection{Preprocessing and cleaning}

We work directly with the processed beta matrix distributed through GEO rather
than reconstructing beta values from raw intensity files. Accordingly, the main
preprocessing steps concern feature filtering and group construction.

For the local-block micro variant, we first keep only autosomal CpGs with valid
genomic coordinates in the Illumina 450k annotation. We drop non-\texttt{cg}
rows, probes with missing genomic position, sex-chromosome probes, rows with
non-finite beta values, and zero-variance rows. Feature screening is then
performed in a fully unsupervised way: each retained CpG is ranked by its
sample variance across the \(656\) subjects, and chronological age is not used
at the screening stage. This avoids outcome-driven feature preselection and
ensures that the subsequent comparison is carried out on a grouped methylation
design defined independently of the response.

\subsection{Grouping design: local methylation blocks}

The grouping rule is chosen to reflect local methylation dependence rather than
only broad gene-level annotation. We therefore sort autosomal CpGs by
chromosome and genomic coordinate, and we start a new group
whenever the next probe is more than \(1000\) base pairs away from the previous
probe on the same chromosome. This produces \emph{local methylation blocks}.
When a block contains an annotated gene label, we name it
\texttt{gene@chr:start-end}; otherwise we retain the coordinate-only label.

After block construction, blocks are ranked by the sum of the variances of
their top ten CpGs. We then select a \emph{micro} subset of exactly \(p=200\)
CpGs by first seeding each chosen block with at least three probes and then
filling the remaining slots by variance within the selected blocks. Under this
rule, the final real-data benchmark contains \(6\) groups with size summary
\((10, 14, 16, 25, 30, 105)\). The dominant groups are concentrated in known
methylation-dense regions, especially the MHC region on chromosome \(6\), with
blocks labeled \texttt{HLA-DQA1}, \texttt{HLA-C}, \texttt{HLA-DRB1}, an
unlabeled local block near \texttt{chr6:29.65Mb}, and a large
\texttt{RNF39}-anchored block; the only non-chromosome-6 group in the micro
subset is a compact \texttt{PM20D1} block on chromosome \(1\).

\subsection{Correlation structure}

Figure~\ref{fig:gse40279-local-block-corr} shows the empirical feature
correlation matrix for the \(200\)-CpG local-block micro design. The grouped
structure is visually clear: the matrix exhibits several dark diagonal blocks,
indicating positive dependence among nearby CpGs within the same genomic
segment, while off-diagonal correlations are comparatively weaker. This is
exactly the kind of structure for which a grouped shrinkage prior is meant to
be useful.

The numerical summaries tell the same story. Across all CpG pairs within the
same block, the mean correlation is approximately \(0.184\), whereas the mean
between-block correlation is only \(0.014\). In absolute value, the contrast is
also substantial: the mean within-block absolute correlation is about \(0.250\),
compared with \(0.076\) between blocks. Thus the local-block construction does
not simply impose a formal partition; it produces groups that correspond to a
real empirical correlation pattern in the methylation design matrix.

\begin{figure}[p]
\centering
\includegraphics[width=0.92\textwidth]{../data/real/gse40279_methylation_age/processed/analysis_bundle/gse40279_micro_local_block_gap1000_feature_correlation_heatmap.pdf}
\caption{Empirical feature correlation heatmap for the GSE40279 local-block
micro design. The dataset contains \(656\) samples, \(200\) CpGs, and \(6\)
genomic proximity groups obtained by merging adjacent autosomal probes whose
genomic gap is at most \(1000\)bp. The visible block-diagonal structure shows
that the grouping rule captures genuine local dependence rather than an
arbitrary partition.}
\label{fig:gse40279-local-block-corr}
\end{figure}

\subsection{Real-data benchmark outcome}

We use this local-block micro dataset as a real-data benchmark for held-out age
prediction under repeated random \(80/20\) train/test splits. The primary role
of this experiment is to assess whether the grouped shrinkage methods behave
reasonably on a biologically structured methylation design whose correlation
pattern is induced by genomic proximity rather than by a synthetic covariance
model.

Table~\ref{tab:gse40279-realdata} reports the held-out prediction results on
the completed common-converged repeats. In this real-data benchmark, all four
methods converged on all \(10/10\) train/test splits, so the paired comparison
uses the full repeated-split sample with no convergence-related attrition.

\begin{table}[H]
\centering
\caption{GSE40279 local-block micro real-data benchmark on common-converged repeated splits. Entries are reported as mean \(\pm\) standard error across 10 paired repeats.}
\label{tab:gse40279-realdata}
\begin{tabular}{lccccc}
\toprule
Method & n paired & RMSE test & MAE test & \(R^2\) test & LPD test \\
\midrule
GR-RHS & 10 & \textbf{10.309089 $\pm$ 0.192506} & \textbf{8.186909 $\pm$ 0.128046} & \textbf{0.515040 $\pm$ 0.016848} & \textbf{-3.777973 $\pm$ 0.013380} \\
Grouped Horseshoe+ & 10 & 10.337849 $\pm$ 0.189610 & 8.208446 $\pm$ 0.138453 & 0.512068 $\pm$ 0.017629 & -3.794836 $\pm$ 0.011130 \\
RHS & 10 & 10.457437 $\pm$ 0.196312 & 8.315410 $\pm$ 0.145084 & 0.500544 $\pm$ 0.018878 & -3.778881 $\pm$ 0.028308 \\
OLS & 10 & 11.704426 $\pm$ 0.206337 & 9.177280 $\pm$ 0.155274 & 0.371653 $\pm$ 0.030980 & -4.276215 $\pm$ 0.061055 \\
\bottomrule
\end{tabular}
\end{table}

The real-data ranking is consistent across the main predictive criteria. The
three Bayesian grouped-shrinkage procedures all substantially outperform OLS on
held-out age prediction, and GR-RHS is the top method on mean RMSE, mean MAE,
mean test \(R^2\), and mean test log predictive density. Grouped Horseshoe+ is
very close to GR-RHS, while RHS remains competitive but slightly weaker on
average across the paired splits.

At this micro scale, the gap among the grouped Bayesian methods is modest,
which is not surprising given that there are only six groups and one very large
\texttt{RNF39}-anchored block contains more than half of the selected CpGs.
Accordingly, the main conclusion is not that the local-block micro benchmark
creates a dramatic separation among GR-RHS, RHS, and Grouped Horseshoe+.
Rather, it shows that the grouped Bayesian framework remains stable and
predictively useful when transferred from synthetic grouped designs to a real
methylation dataset with a clear, empirically verifiable local correlation
structure.

We do not emphasize \texttt{GIGG\_MMLE} in this benchmark because its interval
calibration is poor throughout the simulation study. In the current repository,
\texttt{GIGG\_MMLE} follows an MMLE-style empirical-Bayes path aligned to the
upstream reference implementation, so part of the group-level hyperparameter
uncertainty is plugged in rather than fully propagated through the posterior.
Empirically, this produces intervals that are much too narrow and therefore
severely under-covered relative to the nominal level, even when point-prediction
performance is competitive.

\section{Interpretation}

The simulation evidence supports a scope-conditional empirical statement about
GR-RHS.

When the grouped structure is mild and the active groups are not sharply
distinguishable in their internal signal geometry, GR-RHS behaves like a strong
Bayesian competitor but does not necessarily dominate classical procedures such
as the lasso in overall MSE. This is exactly what one should want from a
well-behaved grouped prior: the group layer does not create an artificial gain
when the setting does not call for it.

When the grouped structure is more informative, however, the GR-RHS mechanism
becomes visible. The mechanism experiments show that this visibility is not a
vague descriptive claim. The group-specific slab ceiling produces ordered
group-level retention in the clean separation setting, suppresses a decoy null
group in the ambiguity setting, and disappears when the group gate is removed.
The broader benchmark then shows that the same mechanism translates into better
performance in the settings where group heterogeneity is genuinely important.

For this reason, the simulation section and the theory section should be read as
paired pieces. The theory identifies \(\kappa_g\) as a group-specific retention
ceiling and emphasizes that the model differs from existing grouped shrinkage
priors by placing the group effect on the finite slab. The simulations show that
this interpretation is not merely formal: it generates an empirically distinct
pattern, and that pattern is most useful exactly in the settings for which the
model was designed.

\end{document}
